\documentclass[11pt]{article}

\usepackage[margin=1in]{geometry}
\usepackage{amsmath,amssymb,amsthm}
\usepackage{graphicx}
\usepackage[hidelinks]{hyperref}

\newtheorem{theorem}{Theorem}
\newtheorem{lemma}{Lemma}
\newtheorem{corollary}{Corollary}

\newcommand{\R}{\mathbb{R}}
\newcommand{\SL}{\mathrm{SL}}
\newcommand{\GL}{\mathrm{GL}}
\newcommand{\Mat}{\mathrm{M}}
\newcommand{\End}{\operatorname{End}}
\newcommand{\card}{\operatorname{card}}

\title{The conjugation action of \(\SL_2(\R)\) is Rosenthal representable}
\author{Run \texttt{fa\_banach\_001}, lane 6}
\date{June 18, 2026}

\begin{document}
\maketitle

\section*{Status}

\textbf{Full candidate solution.}  This packet gives a positive answer
to Question 3.20 of Megrelishvili's paper
\emph{Topological group actions by group automorphisms and Banach
representations}, arXiv:2110.01386.

\begin{center}
\includegraphics[width=0.96\linewidth]{figures/open_problem_crop.png}
\end{center}

\section*{Answer}

Yes.  For \(G=\SL_2(\R)\), the conjugation \(G\)-space \(G_c\) is
Rosenthal representable.

The proof below gives a slightly stronger matrix-group statement.  The
extra generality is included because it makes the mechanism transparent:
conjugation is a finite-dimensional linear action once the group is
viewed inside a matrix algebra.

\section*{Proof Intuition}

The source paper proves that finite-dimensional linear actions are tame
after one-point compactification: the enveloping semigroup consists of
partial linear maps, so it has cardinal at most continuum.  Now
\(\SL_2(\R)\) sits as a closed subset of the four-dimensional vector
space \(\Mat_2(\R)\), and conjugation by \(g\in\SL_2(\R)\) is the
linear map \(A\mapsto gAg^{-1}\) on that vector space.

Therefore the one-point compactification of \(\Mat_2(\R)\) is a tame
compact \(\SL_2(\R)\)-space.  The subset
\(\SL_2(\R)\cup\{\infty\}\) is closed and invariant inside it; closed
invariant subsystems of tame compact systems are tame.  This gives a
tame proper compactification of \(G_c\), which is exactly what is needed
for Rosenthal representability.

\section*{Definitions used}

For a compact metric \(G\)-space \(X\), the source paper recalls the
standard equivalence: \(X\) is representable on a Rosenthal Banach space
if and only if \(X\) is tame; equivalently, its enveloping semigroup
\(E(X)\) has cardinal at most \(2^{\aleph_0}\)
\cite[Fact 3.11]{Megrelishvili}.

A \(G\)-space \(X\) is Rosenthal representable if it admits a proper
\(G\)-compactification by a compact metric tame \(G\)-space, using the
source paper's representation definition and compact-metric tame
characterization.

\section*{The finite-dimensional lemma}

\begin{lemma}
Let \(V\) be a finite-dimensional real vector space and let
\(\alpha V=V\cup\{\infty\}\) be its Alexandrov one-point
compactification.  The natural action
\(\GL(V)\curvearrowright \alpha V\) is tame.
\end{lemma}

\begin{proof}
Let \(p\in E(\alpha V)\).  Choose a net \(T_i\in\GL(V)\) converging
pointwise to \(p\) on \(\alpha V\).  Define
\[
D_p=\{v\in V:p(v)\in V\}.
\]
If \(u,v\in D_p\) and \(a,b\in\R\), then \(T_i u\to p(u)\) and
\(T_i v\to p(v)\) in \(V\), hence
\[
T_i(au+bv)=aT_i u+bT_i v \longrightarrow ap(u)+bp(v)\in V.
\]
Thus \(au+bv\in D_p\) and \(p|_{D_p}\) is linear.  Since \(V\) is
finite dimensional, \(D_p\) is a closed linear subspace and
\(p|_{D_p}:D_p\to V\) is continuous.  For \(v\notin D_p\), by
definition \(p(v)=\infty\).

Consequently each element of \(E(\alpha V)\) is determined by a linear
subspace \(D\leq V\) and a linear map \(D\to V\).  There are only
\(2^{\aleph_0}\) such pairs.  Hence
\(\card E(\alpha V)\leq 2^{\aleph_0}\), and by the source paper's
tameness criterion the compact system \(\alpha V\) is tame.
\end{proof}

\begin{lemma}
Let \(X\) be a compact tame \(G\)-space and let \(Y\subset X\) be a
closed \(G\)-invariant subset.  Then \(Y\) is tame.
\end{lemma}

\begin{proof}
For every \(p\in E(X)\), \(p(Y)\subset Y\): indeed \(p\) is a pointwise
limit of translations by elements of \(G\), each of which sends \(Y\)
to itself, and \(Y\) is closed.  Restriction gives a map
\[
E(X)\longrightarrow Y^Y,\qquad p\mapsto p|_Y .
\]
Every element of \(E(Y)\) is obtained this way after passing to a
convergent subnet in the compact space \(X^X\).  Therefore
\(\card E(Y)\leq \card E(X)\leq 2^{\aleph_0}\), so \(Y\) is tame.
\end{proof}

\section*{Main theorem}

\begin{theorem}
Let \(G\leq \GL_m(\R)\) be a matrix group with the topology inherited
from \(\GL_m(\R)\).  Then the conjugation \(G\)-space \(G_c\) is
Rosenthal representable.
\end{theorem}

\begin{proof}
Put \(V=\Mat_m(\R)\).  The formula
\[
g\cdot A=gAg^{-1}\qquad (g\in G,\ A\in V)
\]
defines a continuous finite-dimensional linear action of \(G\) on
\(V\), because \(A\mapsto gAg^{-1}\) is a linear automorphism of the
vector space \(V\).

Let \(\alpha V=V\cup\{\infty\}\) be the one-point compactification and
let
\[
K=\overline{G}^{\,\alpha V}
\]
be the closure of the image of \(G\subset V\).  The inclusion
\(G\hookrightarrow K\) is a topological embedding with dense image, and
the compact space \(K\) is \(G\)-invariant under the extended linear
action.  On the dense copy of \(G\), this action is exactly
conjugation:
\[
g\cdot x=gxg^{-1}.
\]

By the first lemma, \(\alpha V\) is tame for the full linear group
\(\GL(V)\), hence also for the restricted \(G\)-action.  By the second
lemma, the closed invariant subsystem \(K\subset\alpha V\) is tame.
Thus \(K\) is a tame compact metric \(G\)-compactification of \(G_c\).
By the compact-metric tame/Rosenthal equivalence from the source paper,
\(G_c\) is Rosenthal representable.
\end{proof}

\begin{corollary}
For \(G=\SL_2(\R)\), the conjugation \(G\)-space \(G_c\) is Rosenthal
representable.  Hence Question 3.20 of \cite{Megrelishvili} has a
positive answer.
\end{corollary}

\begin{proof}
\(\SL_2(\R)\) is a matrix group in \(\GL_2(\R)\) with its usual
topology, so the theorem applies.  In this case
\(\SL_2(\R)=\{A\in\Mat_2(\R):\det A=1\}\) is closed in \(\Mat_2(\R)\),
so the compactification \(K\) above is simply the one-point
compactification \(\SL_2(\R)\cup\{\infty\}\).
\end{proof}

\begin{corollary}
For every \(n\geq 2\), the conjugation action of \(\SL_n(\R)\) is
Rosenthal representable.
\end{corollary}

\begin{proof}
Apply the theorem to the usual matrix embedding
\(\SL_n(\R)\leq \GL_n(\R)\).
\end{proof}

\section*{Verification notes}

The proof avoids the common false step ``a tame orbit implies the whole
action is tame.''  Instead it embeds the entire conjugation \(G\)-space
in a compact tame linear system.  For \(G=\SL_2(\R)\), no singular
matrix boundary is added, since determinant one is closed in
\(\Mat_2(\R)\).  For nonclosed or open matrix groups such as
\(\GL_m(\R)\), the same theorem uses the closure in
\(\alpha\Mat_m(\R)\); this is still a valid compactification because
the original group embeds densely and equivariantly.

This full packet supersedes the earlier partial packet proving only
that the unipotent conjugacy orbit is Rosenthal representable.

\section*{Novelty and literature check}

Local run indexes were searched for arXiv id \texttt{2110.01386}, the
title, \(\SL_2(\R)\), \(G_c\), conjugation, matrix groups,
finite-dimensional linear action, and Rosenthal representability.  The
only matching durable packet was the earlier partial unipotent-orbit
packet from this lane.

Web searches on June 18, 2026 for exact phrases including
\(\SL_2(\R)\), \(G_c\), conjugation action, Rosenthal representable,
matrix group, and finite-dimensional linear action found the source
paper but no later paper explicitly resolving Question 3.20.  The arXiv
abstract for the source paper itself states that the \(\SL_2(\R)\)
case is unclear \cite{arxivPage}.  Thus this is recorded as a full
candidate solution rather than a literature-status identification.

\section*{Human review recommendation}

Review the two heredity points: first, that the closure of \(G\) inside
\(\alpha\Mat_m(\R)\) is a legitimate \(G\)-compactification of the
conjugation \(G\)-space; second, that tameness descends to closed
invariant subsystems through restriction of enveloping semigroups.  If
both are accepted, the proof gives a positive solution of Question 3.20.

\begin{thebibliography}{9}

\bibitem{Megrelishvili}
M. Megrelishvili,
\emph{Topological group actions by group automorphisms and Banach
representations},
Forum Mathematicum, accepted version; arXiv:2110.01386, 2023.

\bibitem{arxivPage}
arXiv page for M. Megrelishvili,
\emph{Topological group actions by group automorphisms and Banach
representations},
\url{https://arxiv.org/abs/2110.01386}, accessed June 18, 2026.

\end{thebibliography}

\end{document}

